\input{../000_header_year.txt}

\begin{document}


%++++++++++++++++++++++++++++++++++++++++++++++++++++++++++++++++++++++++++++++++
%++++++++++++++++++++++++++++++++++++++++++++++++++++++++++++++++++++++++++++++++
%++++++++++++++++++++++++++++++++++++++++++++++++++++++++++++++++++++++++++++++++
\exheader
%++++++++++++++++++++++++++++++++++++++++++++++++++++++++++++++++++++++++++++++++
%++++++++++++++++++++++++++++++++++++++++++++++++++++++++++++++++++++++++++++++++
%++++++++++++++++++++++++++++++++++++++++++++++++++++++++++++++++++++++++++++++++

%%%%%%%%%%%%%%%%%%%%%%%%%%%%%%%%%%%%%%%%%%%%%
%　オブザーバ
%　MATLAB mファイルでA,B,C行列（２次元）と指定極を指定して
% オブザーバゲインKを求める
%
%　mファイルの使い方
%　　　まず A,B,C行列とpa,pbを定義する
%%%%%%%%%%%%%%%%%%%%%%%%%%%%%%%%%%%%%%%%%%%%%
%================================
%	matlab100.mを利用


\input{105_def_file.txt}

%================================

\def\state{
	\left (
	\begin{array}{c}
	 x_1 \\ x_2 
	\end{array}
	\right )
}

\def\cK{
	\left ( \begin{array}{c}
		k_1 \\ k_2 
	\end{array} \right )
}

\def\cKT{
	\left ( \begin{array}{cc}
		k_1 & k_2 
	\end{array} \right )
}

%================================
\noindent
{\bf オブザーバ（演習）}　次のシステムに対して，オブザーバを設計せよ．ただし，オブザーバの極は$\chrpa,\chrpb$とせよ．
\begin{align*}
	\frac{d}{dt}
	\state
	& = 
	\chrA \state
	+
	\chrB u
\\
	y&=
	\chrC \state
\end{align*}

\vfill

\begin{tabular}{|l|}
\hline
解答：設計したオブザーバ\\[1.5cm]
\hspace*{15cm}\\
\hline
\end{tabular}

%---------------------------------------------------------------------------
\newpage
\noindent
{\bf オブザーバ（演習）}　次のシステムに対して，オブザーバを設計せよ．ただし，オブザーバの極は$\chrpa,\chrpb$とせよ．
\begin{align*}
	\frac{d}{dt}
	\state
	& = 
	\chrA \state
	+
	\chrB u
\\
	y&=
	\chrC \state
\end{align*}

$A+KC$の極配置をするために次の仮想システムを考える．
\begin{align*}
 \frac{d \xi}{dt} & = A^T \xi + C^T \eta, 
\hspace{1cm}
	A^T = \chrAT
	, \ 
	C^T = \chrCT
\end{align*}

この仮想システムの閉ループ極を$\chrpa$, $\chrpb$とする仮想フィードバックを設計
\begin{align*}
 \eta =  K^T \xi
	= \cKT \xi
\end{align*}

仮想システムの閉ループ系
\begin{align*}
 \frac{d \xi}{dt} & = A^T \xi + C^T K^T \xi
\\
	&= \chrAT \xi
	+
	\chrCT \cKT \xi
\\
	&= \chrAT \xi + \chrCKT \xi
\\
	&= \chrACKT \xi
\end{align*}

仮想システムの閉ループ系の特性多項式
\begin{align*}
	&\det
	\left \{
		sI- \chrACKT
	\right \}
	=
	\det \chrsIACKT
\\
	& \hspace{1cm} =
	(\chrsIACKTa) \cdot (\chrsIACKTd) -(\chrsIACKTb)(\chrsIACKTc)
	=
	\chrdetsIACKT
\end{align*}

極が$\chrpa$, $\chrpb$ となる特性多項式は
\begin{align*}
	(s+\chrmpa)(s+\chrmpb) = s^2 + \chrcoefbd s + \chrcoefcd
\end{align*}

係数を比較して
\begin{align*}
	\chrcoefb=\chrcoefbd, \ \chrcoefc=\chrcoefcd 
	\ \ \Rightarrow \ \ 
	k_1=\chrsKa, \ k_2=\chrsKb
\end{align*}

よって
\begin{align*}
	K&= \cK = \chrsK
\end{align*}

オブザーバは
\begin{align*}
 \frac{d \hat{x}}{dt} &= A \hat{x} + B u - K (y - C \hat{x})
\\
 \frac{d}{dt}
	\left (
	\begin{array}{c}
	 \hat{x}_1 \\ \hat{x}_2 
	\end{array}
	\right )
	&=
	\chrA \state
	+ 
	\chrB u
	-
	\chrsK
	\left \{
		y -
			\chrC
			\left (
			\begin{array}{c}
			 \hat{x}_1 \\ \hat{x}_2 
			\end{array}
			\right )
	\right \}
\end{align*}




\end{document}

